\ifdefined\NoetherPaperTwentyOneBodyOnly
\else
\documentclass[11pt]{article}
\usepackage[a4paper,margin=1.45cm]{geometry}
\usepackage[T1]{fontenc}
\usepackage[utf8]{inputenc}
\usepackage{lmodern}
\usepackage[french]{babel}
\usepackage{amsmath,amssymb,mathtools}
\usepackage[protrusion=true,expansion=false]{microtype}
\setlength{\parindent}{1.25em}
\setlength{\parskip}{0.10em}
\allowdisplaybreaks
\setlength{\emergencystretch}{12em}
\sloppy\hbadness=10000\vbadness=10000\hfuzz=2pt\vfuzz=10pt\raggedbottom
\makeatletter\renewcommand\footnotesize{\@setfontsize\footnotesize{7.5}{8.7}\selectfont}\makeatother

\begin{document}
\fi

% Unit: Paper 21. Direct French translation from the German final-audited source slice.

\section*{21. Calcul des variations formel et invariants différentiels}
\begin{center}
Encyklopädie d. math. Wiss. II, 3 (1922), p.~68--71 (dans: R. Weitzenböck, \emph{Differentialinvarianten})
\end{center}

\setcounter{footnote}{148}
\subsection*{28. Calcul des variations formel et invariants différentiels.\footnote{Ce numéro est dû à E. Noether.}}

La production de tous les invariants différentiels et la déduction des théorèmes principaux peuvent être menées de la manière la plus uniforme par les méthodes formelles du calcul des variations. Ce principe remonte à Riemann\textsuperscript{78)} et a reçu son développement principal dans la théorie des formes par Lipschitz\textsuperscript{73)}.

Si l'on part du problème variationnel appartenant à une forme homogène $f(dx)$ -- où l'on suppose
\[
\left|\frac{\partial^2 f}{\partial dx_i\,\partial dx_k}\right|\ne0
\]
-- à savoir
\begin{equation}
 \delta\int_{t_0}^{t_1} f(x')\,dt=0,
 \qquad \left(x_i'=\frac{dx_i}{dt}\right),
\tag{140}
\end{equation}
on est conduit, par intégration partielle, au système invariant d'équations
\[
 2\psi_i=\frac{d}{dt}\frac{\partial f}{\partial x_i'}-\frac{\partial f}{\partial x_i}=0,
\]
où les expressions de Lagrange $\psi_i$ sont contragrédientes aux $dx$. Formellement, cela se voit le plus simplement si l'on définit les $\psi_i$ par l'identité formelle sans intégrale correspondant à l'intégration partielle (équation centrale de Lagrange\footnote{La désignation, pour un $f$ spécial, est due à Heun; voir son article IV, 1 II, p.~447.})
\[
 \delta f'-df_\delta=-2\sum\psi_i(d,d)\,\delta x_i,
 \qquad
 \left(f_\delta=\sum\frac{\partial f}{\partial dx_i}\,\delta x_i\right),
\]
où $f$ aussi bien que $df_\delta$ sont des invariants, et par conséquent aussi le membre de droite. En particulier, pour une forme quadratique, il vient
\[
 \delta\sum g_{ik}dx_i dx_k-2d\sum g_{ik}dx_i\delta x_k
 =-2\sum\psi_\mu(d,d)\,\delta x_\mu.
\]
Si l'on remplace ici $d$ par le cogrédient $(d+\lambda\delta)$ et que l'on compare les puissances de $\lambda$, on obtient le système correspondant d'identités:
\begin{equation}
\left\{
\begin{aligned}
D\sum g_{ik}dx_i dx_k-2d\sum g_{ik}dx_iDx_k
   &=-2\sum\psi_\mu(d,d)\,Dx_\mu,\\
D\sum g_{ik}dx_i\delta x_k-\delta\sum g_{ik}dx_iDx_k
   -d\sum g_{ik}\delta x_iDx_k
   &=-2\sum\psi_\mu(d,\delta)\,Dx_\mu,\\
D\sum g_{ik}\delta x_i\delta x_k-2\delta\sum g_{ik}\delta x_iDx_k
   &=-2\sum\psi_\mu(\delta,\delta)\,Dx_\mu.
\end{aligned}
\right.
\tag{141}
\end{equation}
Il en résulte que $\psi_\mu(d,\delta)$, donc en particulier aussi $\psi_\mu(d,d)$ et $\psi_\mu(\delta,\delta)$, se présentent comme des vecteurs contragrédients. De (141) on calcule
\[
 \psi_\mu(d,\delta)=\sum g_{i\mu}\,d\delta x_i+
 \sum\left[\begin{smallmatrix}ik\\ \mu\end{smallmatrix}\right]dx_i\delta x_k,
\]
et il en résulte le vecteur cogrédient:
\begin{equation}
 p^\sigma(d,\delta)=\sum g^{\sigma\mu}\psi_\mu
 =d\delta x_\sigma+
 \sum\left\{\begin{smallmatrix}ik\\ \sigma\end{smallmatrix}\right\}dx_i\delta x_k.
\tag{142}
\end{equation}
Poser $p^\nu(d,d)=0$ donne évidemment les équations des lignes géodésiques.

La connaissance du vecteur cogrédient $p^\sigma(d,\delta)$ conduit directement à la dérivée covariante (cf. n\textsuperscript{o}~19). Si, par exemple, $h(d,\delta)$ est une forme des deux rangées $dx$ et $\delta x$, alors l'expression
\begin{equation}
 h^{(1)}(dx,\delta x)=\delta h-
 \sum\frac{\partial h}{\partial dx_\sigma}p^\sigma(d,\delta)-
 \sum\frac{\partial h}{\partial \delta x_\sigma}p^\sigma(\delta,\delta)
\tag{143}
\end{equation}
est elle-même un invariant, comme différence d'invariants, et elle est formée de telle sorte que les différentielles secondes s'éliminent. Ainsi $h^{(1)}(dx,\delta x)$ représente une forme ne contenant que des différentielles premières, la dérivée covariante de $h$. Il en va de même lorsque $h$ contient davantage de rangées $d^{(1)}x,\ldots,d^{(\lambda)}x$.\footnote{Lipschitz, J. f. Math. 72 (1870), p.~17.}

La formation de la forme de courbure chez Riemann et Lipschitz repose sur le même principe. On passe de $\delta^2 f$ à la \og forme normale de la seconde variation \fg{}
\begin{equation}
 \Omega=\delta^2\sum g_{ik}dx_i dx_k
 -2d\delta\sum g_{ik}dx_i\delta x_k
 +d^2\sum g_{ik}\delta x_i\delta x_k,
\tag{144}
\end{equation}
qui, comme agrégat d'invariants, représente elle-même encore un invariant, et dans laquelle, maintenant par généralisation de l'équation centrale de Lagrange, les différentielles troisièmes sont éliminées. L'élimination des différentielles secondes se fait par analogie avec la dérivation covariante, en formant la différence
\begin{equation}
 K=\Omega-2\sum\{p^\sigma(d,\delta)\psi_\sigma(d,\delta)-p^\sigma(\delta,\delta)\psi_\sigma(d,d)\},
\tag{145}
\end{equation}
ce qui peut aussi s'écrire\footnote{Lipschitz, J. f. Math. 82 (1876), \S~1.}
\begin{equation}
\begin{aligned}
 \frac12 K={}&d\sum\psi_\sigma(d,\delta)\delta x_\sigma
 -\delta\sum\psi_\sigma(d,d)\delta x_\sigma\\
 &-\sum\{p^\sigma(d,\delta)\psi_\sigma(d,\delta)-p^\sigma(\delta,\delta)\psi_\sigma(d,d)\}.
\end{aligned}
\tag{146}
\end{equation}
La loi de formation de $K$ peut encore s'énoncer en disant que, dans $\Omega$, les différentielles secondes sont éliminées par l'annulation de $p(d,\delta)$ et de $p(\delta,\delta)$ -- respectivement des membres de droite de (141). C'est la définition de la forme de courbure par Riemann (cf. n\textsuperscript{o}~21), mais il faut remarquer expressément -- ce que (146) montre avec une clarté particulière -- que cette annulation ne peut avoir lieu qu'après la différentiation; le mode de formation formel de Lipschitz contourne cette difficulté. De même, la dérivée covariante $h^{(1)}$ (143) pourrait être définie en éliminant dans $\delta h$ les différentielles secondes par l'annulation de $p(d,\delta),\ldots$.

De la forme de courbure résulte, par formation successive des dérivées covariantes, une suite de formes $[\Phi_4,\Phi_5,\ldots$ chez Christoffel], dont les invariants projectifs, après adjonction de $f$, épuisent selon Christoffel et Ricci la totalité des invariants différentiels de la forme quadratique, et dont l'équivalence vis-à-vis des transformations linéaires, sans autre condition d'intégrabilité, exprime l'équivalence de deux formes différentielles quadratiques (cf. n\textsuperscript{o}~22). La raison interne de ce théorème de réduction repose sur l'existence des coordonnées normales de Riemann issues du problème variationnel (140) (cf. n\textsuperscript{o}~20), lesquelles transforment en droites les lignes géodésiques passant par un point, et se transforment par conséquent linéairement sous une transformation arbitraire des $x$. La formation de tous les invariants et la question de l'équivalence sont donc ramenées à la question correspondante pour les composantes homogènes individuelles dans le développement de $f$ suivant les coordonnées normales, et l'on montre que ces composantes se composent linéairement à partir de $f,\Phi_4,\Phi_5,\ldots$. Pour les formes quadratiques, cela a été exécuté en détail par Vermeil\textsuperscript{92)}, en s'appuyant sur une note de E. Noether\textsuperscript{93)}. D'après cette note, les théorèmes et les méthodes subsistent lorsque l'on prend pour base une expression différentielle arbitraire (cf. n\textsuperscript{os}~22 et 23); ainsi apparaît la portée des méthodes du calcul des variations formel par rapport à celles de la pure théorie de l'élimination.\footnote{Comparer, pour ces développements et pour d'autres interprétations géométriques, les conférences de séminaire de F. Klein indiquées sous \og Littérature \fg{}.}

L'annulation de $p(d,\delta)$ a été interprétée géométriquement par Levi-Civita\textsuperscript{134)} comme \og déplacement parallèle d'un vecteur \fg{} (cf. aussi Hessenberg\textsuperscript{133)}). Cette notion, que Weyl a saisie axiomatiquement, est essentiellement restreinte aux formes quadratiques, mais elle admet une généralisation d'un autre genre. Elle subsiste lorsqu'on élargit le groupe (multiplication des $g_{ik}$ par une fonction arbitraire), dès que l'on prend pour base, outre la forme quadratique, encore une forme linéaire. $p(d,\delta)$ est alors déterminé de manière unique, et l'on peut effectuer des formations d'invariants correspondant à celles qui précèdent et définir des coordonnées géodésiques; mais $p(d,d)=0$ ne provient alors plus d'un problème variationnel!\footnote{H. Weyl, \emph{Reine Infinitesimalgeometrie}, Math. Ztschr. 2 (1918), p.~384--411, et \emph{Raum, Zeit, Materie} (3\textsuperscript{e} et 4\textsuperscript{e} éd.).} Les questions concernant le théorème de réduction et l'équivalence ne sont pas encore abordées ici de manière générale; R. Weitzenböck n'a donné un théorème de réduction que pour les invariants d'ordre le plus bas.\footnote{Wien. Ber. 129 (1920), p.~683 et 697.}

Enfin, il faut signaler les \og problèmes variationnels invariants \fg{} généraux, c'est-à-dire ceux dans lesquels l'intégrale $J$ (simple ou multiple) est invariante vis-à-vis d'un groupe quelconque au sens de Lie. Des cas spéciaux, intéressants pour la physique, font l'objet de travaux de Hamel\footnote{Math. Ann. 59 (1904), p.~416--434; Ztschr. Math. Phys. 50 (1904), p.~1--54.}, Herglotz\footnote{Ann. d. Phys. (4) 36 (1911), p.~511, en particulier \S~9.}, Lorentz\footnote{Verslag der Amsterd. Akad., avril, sept. 1916, oct. 1916, mai 1917.}, Fokker\footnote{Ibid., janv. 1917.}, Weyl\footnote{Ann. d. Phys. (4) 54 (1917), p.~117, \S~2.} et Klein\footnote{Gött. Nachr. 1917, p.~469--482; ibid. 1918, p.~171--189.}. La formulation principielle chez E. Noether\footnote{Gött. Nachr. 1918, p.~235--257.} montre qu'à l'invariance de $J$ vis-à-vis d'un $G_\varrho$ (groupe fini à $\varrho$ paramètres essentiels) correspondent $\varrho$ divergences linéairement indépendantes; à l'invariance vis-à-vis d'un groupe infini qui contient $\varrho$ fonctions arbitraires jusqu'à la $\sigma$-ième dérivée correspondent $\varrho$ dépendances entre les expressions de Lagrange et leurs dérivées jusqu'à l'ordre $\sigma$. Dans les deux cas, la réciproque vaut. Comme les expressions de Lagrange deviennent des invariants relatifs du groupe, on possède en même temps un procédé générateur d'invariants.

\begin{center}
(Achevé en mars 1921.)
\end{center}

\clearpage

\ifdefined\NoetherPaperTwentyOneBodyOnly
\else
\end{document}
\fi
